\subsection{System Feature 3: Mentee RSVP/un-RSVP} 


\subsubsection{Description}
Mentees can RSVP and un-RSVP to events created by mentors and time conflicts
will be handled.
\textit{(Priority: High)}



\subsubsection{Functional Requirements}
\begin{itemize}[label=$\bullet$]
    \item \textbf{RSVP-1:} Mentees shall have the ability to RSVP to an event
    that they are eligible for.
    
    \item \textbf{RSVP-2:} Mentees shall have the ability to un-RSVP from an
    event.
    
    \item \textbf{RSVP-3:} Error message shall be displayed when a mentee
    attempts to RSVP an event whose time slot (start time to end time) overlaps
    with a RSVP'd events's time slot.
    
    \item \textbf{RSVP-4:} Mentors shall be able to approve a mentee's
    registration in a time-conflicting event.
\end{itemize}



\subsubsection{Use Cases}

% Use Case 6
\begin{usecase}{UC-6: RSVP Event (with Time Conflict Handling)} 
    \textbf{Primary Actor}  & Mentee \newline \textbf{Secondary Actor:} Mentor   \\ \hline
    \textbf{Description}    & Mentees can RSVP to events that they are eligible
                              for. A mentee may RSVP for an event that has time
                              conflict with one of their currently RSVP'd events
                              with mentor approval. The mentor will receive a
                              notification on their mentor dashboard that allows
                              them to approve or reject the mentee's RSVP
                              request for the overlapping event.                 \\ \hline
    \textbf{Preconditions}  & PRE-1: Mentee is elegible for this event (event is
                              displayed on their personalized calendar).         \\ \hline
    \textbf{Postconditions} & POST-1: Mentee is added to event's RSVP list.      \newline 
                              POST-2: Event is marked on mentee's personalized
                              calendar as RSVP'd.                                \\ \hline
    \textbf{%
        Normal Flow
        (No Conflicts)%
    }                       & 
    \begin{tabenum}
        \item Mentee selects an eligible event on their personalized calendar.
        \item Mentee clicks on "RSVP".
        \item Mentee gets added to event RSVP list.
        \item Application marks event as RSVP'd on mentee's personal calendar.
    \end{tabenum}                                                                \\ \hline
    \textbf{%
        Alternative Flow
        (Mentor Approves)%
    }                       & 
    \begin{tabenum}
        \item Mentee clicks on eligible event on their personalized calendar.
        \item Application displays time conflict warning and registration is
        blocked.
        \item Mentee clicks on "Request Mentor Approval" button.
        \item Notification is sent to the mentor's dashboard.
        \item The mentor approves the registration.
        \item Mentee is added to the new event's RSVP list in the database.
        \item Application marks new event as RSVP'd on mentee's personal
    calendar.
    \end{tabenum}                                                                \\ \hline
    \textbf{%
        Alternative Flow
        (Mentor Rejects)%
    }                       & 
    \begin{tabenum}
        \item Mentee clicks on eligible event on their personalized calendar.
        \item Application displays time conflict warning and registration is
        blocked.
        \item Mentee clicks on "Request Mentor Approval" button.
        \item Notification is sent to mentor's dashboard.
        \item Mentor rejects the registration.
        \item Application sends rejection notification to mentee and no changes
        are made to the mentee's personalized calendar.
    \end{tabenum} \\
\end{usecase}

% Use Case 7
\begin{usecase}{UC-7: Un-RSVP Event} 
    \textbf{Primary Actor}  & Mentee                                             \\ \hline
    \textbf{Description}    & Mentees can un-RSVP to any currently RSVP'd
                              events. The mentees' personalized calendar must
                              reflect these changes.                             \\ \hline
    \textbf{Preconditions}  & PRE-1: Mentee has already RSVP'd to an event.      \\ \hline
    \textbf{Postconditions} & POST-1: Event is marked as eligible on the
                              mentee's personalized calendar.                    \newline
                              POST-2: Mentee is removed from event's RSVP list.  \\ \hline
    \textbf{Normal Flow}    & 
    \begin{tabenum}
        \item Mentee selects event that they have already RSVP'd.
        \item Mentee selects "Un-RSVP" button.
        \item Event is marked as eligible on mentee's personalized calendar.
        \item Mentee is removed from event's RSVP list in database.
    \end{tabenum} \\
\end{usecase}



\subsubsection{User Stories}

\begin{itemize}[leftmargin=*, label=$\square$]
    \item \textbf{As a mentee}, I want to RSVP events so that I can secure my
    spot.
    \begin{itemize}[label=$\circ$, nosep]
        \item \textit{Acceptance:} A mentee can RSVP to an event that is marked
        as eligible on their personalized calendar.

        \item \textit{Acceptance:} A mentee is saved to event's RSVP list in
        database.

        \item \textit{Acceptance:} After mentee RSVPs, event is vistually
        distinguishable from other eligible events on their personal calendar.
    \end{itemize}

    \item \textbf{As a mentee}, I want to un-RSVP from events, so that my mentor
    knows I no longer want to attend.
    \begin{itemize}[label=$\circ$, nosep]
        \item \textit{Acceptance:} A mentee can un-RSVP from an event.
        
        \item \textit{Acceptance:} A mentee is removed from event's RSVP list in
        database.

        \item \textit{Acceptance:} After mentee un-RSVPs, their personal
        calendar should mark the event as eligible.
    \end{itemize}

    \item \textbf{As a mentee}, I want to request mentor approval for an
    overlapping event so that I can possibly attend both.
    \begin{itemize}[label=$\circ$, nosep]
        \item \textit{Acceptance:} A mentee can request permission to override
        time conflict after clicking on event details from their personalized
        calendar.

        \item \textit{Acceptance:} Application marks time-conflicting event as
        RSVP'd if mentor approves their RSVP request.
    \end{itemize}

    \item \textbf{As a mentor}, I want to receive time conflict notifications,
    so that I can approve or reject a mentee's RSVP request for an event with a
    time conflict.
    \begin{itemize}[label=$\circ$, nosep]
        \item \textit{Acceptance:} A mentor can receive a time conflict
        notification when a mentee attempts to RSVP for an event that overlaps
        with their RSVP'd events.

        \item \textit{Acceptance:} A mentor can approve or reject a mentee's
        RSVP request for a time conflicting event.
    \end{itemize}
\end{itemize}
